\section{GPU 计算架构与线程块驻留}

\subsection{流式多处理器是主要硬件组织单元}

一块 GPU 由多个\term{流式多处理器}{Streaming Multiprocessor, SM}组成。每个流式多处理器内部包含若干执行核心、控制资源和片上存储资源；多个流式多处理器通过更高层次的存储系统访问全局内存 (global memory)。以 Hopper H100 的架构示意为例，144 个流式多处理器、每个流式多处理器 128 个核心，对应
\[
144\times 128=18\,432
\]
个核心。

\begin{figure}[H]
  \centering
  \includegraphics[width=0.88\textwidth]{gpu_architecture_p06.png}
  \caption{多个流式多处理器共享更高层次的全局内存；每个流式多处理器内部具有控制、执行核心和片上存储资源}
\end{figure}

\begin{definitionbox}
\term{流式处理器}{Streaming Processor, SP}或 CUDA 核心表示流式多处理器内部的标量执行核心。流式多处理器负责维护大量驻留线程的执行状态并调度线程束；流式处理器是实际完成相应算术工作的执行资源之一。两者不能用同一个“核心数”概念替代。
\end{definitionbox}

\subsection{线程块以整体形式分配到一个流式多处理器}

CUDA 线程不是逐个、独立地被分配到任意流式多处理器。硬件以\textbf{线程块粒度}完成驻留分配：一个线程块的全部线程都属于同一个流式多处理器，并在该流式多处理器上保持驻留直到线程块完成。一个线程块不能拆到多个流式多处理器上。

多个线程块可以同时驻留在同一个流式多处理器上，但前提是它们共同占用的硬件资源没有超过该流式多处理器的限制。尚未获得驻留资源的线程块处于等待状态；当已有线程块完成并释放资源后，等待线程块才可能获得驻留机会。

\begin{figure}[H]
  \centering
  \includegraphics[width=0.92\textwidth]{block_assignment_p07.png}
  \caption{线程块以整体形式分配到流式多处理器；超出当前驻留能力的线程块必须等待资源释放}
\end{figure}

这里必须区分\textbf{驻留}和\textbf{正在执行}。假设某个流式多处理器能够容纳 2048 个驻留线程，并不意味着这 2048 个线程在同一个时钟周期中都各自占用一个物理执行核心。一个流式多处理器的物理执行核心数通常远小于可驻留线程数；驻留线程会进一步组成线程束，由调度器选择其中已经具备执行条件的线程束发射指令。

\begin{keybox}
“线程块被分配到某个流式多处理器”表示该线程块所需的执行状态和资源已在该流式多处理器上建立，并不表示线程块中的所有线程同时占用独立核心执行。
\end{keybox}

\subsection{为什么同一线程块必须留在同一个流式多处理器}

同一线程块中的线程依赖两类低开销协作资源：

\begin{itemize}
  \item \term{共享内存}{shared memory}：物理资源位于流式多处理器上，但逻辑上按线程块分配；同一线程块中的线程可以访问该线程块的共享内存。
  \item 线程块级屏障：同一线程块的线程可以使用 \code{__syncthreads()} 等机制等待彼此达到同步点。
\end{itemize}

如果一个线程块被拆分到不同流式多处理器，这些低开销机制就失去统一的本地硬件基础。因此，线程块边界既是编程模型中的协作边界，也是硬件资源分配的重要边界。

不同线程块之间如果需要交换数据，通常必须依赖全局可见的存储与更高层次的同步机制。其成本和约束都显著高于线程块内部协作，因此 CUDA 的基本执行模型要求普通线程块之间不依赖具体的调度顺序。

\subsection{驻留需要同时满足多类资源限制}

线程块获得驻留位置之前，流式多处理器必须能够为整个线程块预留足够资源。主要约束包括：

\begin{itemize}
  \item 每线程块最大线程数；
  \item 每流式多处理器最大驻留线程数；
  \item 每流式多处理器最大驻留线程块数；
  \item 寄存器文件容量；
  \item 共享内存容量；
  \item 最大驻留线程束数等架构限制。
\end{itemize}

因此，“线程数还放得下”只是必要条件之一。例如，一个线程块可能在线程数方面合法，但其寄存器或共享内存需求使得同一个流式多处理器只能驻留很少的线程块。后面的占用率分析将把这些约束统一起来。
